\documentclass[10pt,a4paper]{amsart}
\usepackage[margin=19mm]{geometry}
\usepackage{amsmath,amssymb,mathtools}
\usepackage[scr=rsfs,cal=euler]{mathalfa}
\usepackage{xfrac}
\usepackage[unicode,hidelinks,bookmarks=false]{hyperref}
\newcommand{\ie}{i.e.\ }
\newcommand{\cf}{cf.\ }
\newcommand{\resp}{respectively\ }
\newcommand{\Subsection}{\subsection}
\providecommand{\nobreakdash}{\nobreak\hbox{-}}
\setlength{\emergencystretch}{2em}
\multlinegap0pt
\usepackage{tikz}
\usepackage{amssymb}
\usepackage{xcolor}
\usepackage[normalem]{ulem}
\newtheorem{corollary}{Corollary}
\newtheorem{lemma}{Lemma}
\DeclareMathOperator*{\Res}{Res}
\title{Laplace measure transitions and ghosts\\ for meromorphic functions}
\author{Jo\~ao Fontinha, Jorge Buescu, Jaouen Ramalho}
\date{Source: arXiv:2604.04162v2 (CC BY 4.0)}
\begin{document}
\maketitle

\noindent\emph{Source and licence.} Laplace measure transitions and ghosts\\ for meromorphic functions. Version arXiv:2604.04162v2 (CC BY 4.0), distributed by arXiv under the Creative Commons Attribution 4.0 International licence, http://creativecommons.org/licenses/by/4.0/, as recorded in the arXiv licence metadata for this exact version and shown on the abstract page https://arxiv.org/abs/2604.04162v2. Full licence notice: \texttt{licenses/arxiv-2604.04162-CC-BY-4.0.txt}.\par\smallskip

\noindent\textit{Editorial scope.} Counted unit: the article's lemma lem:transitionforfinite (source0.tex lines 340--348). The article's complete original proof of it is delivered with it (source0.tex lines 350--418, one \texttt{proof} environment of 69 lines). No mathematics is written, repaired or re-derived: the statement and the proof are the article's own lines, in the article's own notation, and no retained line is substituted or re-flowed.  Retained uncounted as the source-local context the proof needs: lines 183--186; lines 300--311.  Nothing was omitted from any retained span, so the package carries no omission record.  No retained line contains a citation, so the wrapper carries no bibliography and no bibliography entry.  No retained line contains a figure environment, an image-inclusion command or drawing code, so no image is delivered and the package declares no asset dependency.  Editorial formatting only, in addition: an AMS \texttt{amsart} preamble that restates the article's own declarations and macros used by the retained lines, namely the environments \texttt{corollary}, \texttt{lemma} on separate counters, together with the standard packages those lines need.  The retained environments print in this document as Corollary 1, Lemma 1 in order of appearance, whereas the article prints the same items with the numbering of its own full document.  The complete native source of the article as distributed in the arXiv e-print for this exact version is bundled unchanged as \texttt{sources/197864/source0.tex}.
     \begin{equation}
     \label{eq_PL}
         |F(x+iy)| \leq Ae^{Be^{K|y|}} 
     \end{equation}

\begin{corollary}\label{cor:finitepoles}
Let \( F(s) \) be a meromorphic Laplace pair \( (\mu_l,\mu_r) \) having finitely many poles  on the imaginary axis 
\( s=ip_n \), \( 1\le n\le N \), and satisfying the Phragmén--Lindelöf condition \eqref{eq_PL}. Then for each real \( t \)
\begin{equation}
    \label{eq:finite_transition_formula}
\mu_r(t)-\mu_l(t) 
=
\sum_{n=1}^{N}
\Res\!\left((e^{st}-1)\frac{F(s)}{s},\, s=ip_n\right).
\end{equation}

\end{corollary}

\begin{lemma}\label{lem:transitionforfinite}
Let $F(s)$ be a meromorphic Laplace pair $(f_l,f_r)$ satisfying the hypotheses of Corollary~\ref{cor:finitepoles}. 
For each pole $s=ip_n$, let $r_n$ denote its order and let
$\displaystyle{P_n(s)= \sum_{k=1}^{r_n}\frac{a_{n,k}}{(s-ip_n)^k}}$
be the principal part of its Laurent expansion around $s=ip_n$. Then
\begin{equation}\label{eq:transitionformula}
f_r(t)-f_l(t)= \sum_{n=1}^{N}e^{ip_nt}\sum_{k=1}^{r_n}a_{n,k}\frac{t^{k-1}}{(k-1)!}.
\end{equation}
\end{lemma}

\begin{proof}
Since $\mu_{r,l}(t)= \int_{0}^{t} f_{r,l}(u)\,du$, Corollary~\ref{cor:finitepoles} yields
\[
f_r(t)-f_l(t)
=
\frac{d}{dt}
\sum_{n=1}^{N}
\Res\!\left(e^{st}\frac{F(s)}{s};\, s=ip_n\right).
\] We proceed with the construction of an explicit expression for 
$\displaystyle{
\Res\!\left(e^{st}\frac{F(s)}{s};\, s=ip_n\right).}$ Expanding the functions $e^{st}$, $1/s$, and $F(s)$ in Laurent series around $s=ip_n$, we obtain
\begin{equation*}
    \label{eq_finite_transition}
e^{st}\frac{F(s)}{s}
=
e^{itp_n}
\left[
\sum_{\substack{l,m\ge 0 \\ k\ge -r_n}}
\frac{(-1)^m}{(ip_n)^{m+1}}
\frac{t^l}{l!}
b_{n,k}
(s-ip_n)^{m+l+k}
\right],
\end{equation*}


\noindent where $b_{n,k}=a_{n,-k}$ and $-r_n\le k\le -1.$ 

The residue at $s=ip_n$ is the coefficient of $(s-ip_n)^{-1}$, and so
\[
\Res\!\left(e^{st}\frac{F(s)}{s},\, s=ip_n\right)
=
e^{itp_n}
\left[
\sum_{\substack{l,m\ge 0 \\ k\ge -r_n \\ m+l+k=-1}}
\frac{(-1)^m}{(ip_n)^{m+1}}
\frac{t^l}{l!}
b_{n,k}
\right].
\]

\noindent Since $l$ and $m$ are non-negative, the condition $m+l+k=-1$ implies
$0\le l\le -k-1$, and therefore it follows that
\begin{equation}\label{eq:residuecomputation}
\Res\!\left(e^{st}\frac{F(s)}{s},\, s=ip_n\right)
=
e^{itp_n}
\left[
\sum_{k=-r_n}^{-1}
\sum_{l=0}^{-k-1}
b_{n,k}
(-1)^{k+l+1}
(ip_n)^{k+l}
\frac{t^l}{l!}
\right].
\end{equation}

\noindent Differentiating with respect to $t$ we obtain
\begin{equation}\label{eq:differentiationofresidues}
\frac{d}{dt}\Res\!\left(e^{st}\frac{F(s)}{s},\, s=ip_n\right)= e^{itp_n}\sum_{k=-r_n}^{-1}
b_{n,k}
\frac{t^{-k-1}}{(-k-1)!}
=
\sum_{k=1}^{r_n}
a_{n,k}
\frac{t^{k-1}}{(k-1)!}.
\end{equation}
Finally, summing over all poles yields \eqref{eq:transitionformula}, completing the proof
\end{proof}

\end{document}
